\section{The truncated twist and the logarithmic derivative at one}
\label{sec:window-explicit}

Let \(\chi\) be a nonprincipal Dirichlet character of modulus \(m\ge1\).
Write \(\ell_\chi(s)=-L'(s,\chi)/L(s,\chi)\) wherever defined,
and \(a_n=\Lambda(n)\chi(n)\). Write \(Q\ge2\) for the height
and conductor cutoff. The box hypothesis is
\[
 m\le Q,\qquad L(s,\chi)\ne0
 \quad\bigl(\Re s>399/400,\ |\Im s|\le Q\bigr).
\]
For \(v>0\) define
\[
 A_\chi(v)=\sum_{n\le v}a_n(v/n-1),\qquad
 G_\chi(v)=A_\chi(v)-v\ell_\chi(1),\qquad
 R_\chi(z)=\ell_\chi(1)-\sum_{n\le z}a_n/n .
\]
Positive-integer summation includes the floor of each real cutoff.
For \(v\le z\), direct separation at \(n=v\) gives the exact identity
\begin{equation}
 G_\chi(v)=-\sum_{n\le z}a_n\min(1,v/n)-vR_\chi(z).
 \label{eq:box-finite-split}
\end{equation}
The choice of sign records that the source centers by adding
\(vL'(1,\chi)/L(1,\chi)\).

\subsection{An explicit contour estimate}

For \(0<\delta\le1/800\), \(z\ge1\), \(0<\eta\le1\), and
\(h=z\eta\), we claim
\begin{equation}
 |A_\chi(z+h)-A_\chi(z)-h\ell_\chi(1)|
 \le C_\delta\log Q
   \left\{(\log Q)\eta z^{399/400+\delta}+z^2/Q\right\}.
 \label{eq:box-contour}
\end{equation}
Put \(\sigma=399/400+\delta\), \(T=Q/2\) and \(L_Q=\log Q\).
First fix the bound
\[
 |\ell_\chi(w)|\le B_\delta L_Q\quad
          (\Re w\ge \sigma,\ |\Im w|\le Q-1).
\]
In the range \(\Re w<11/10\), the source's local logarithmic
derivative lemma for \(L\) at \(w\) expands \(L'/L\) as a sum of
terms \(m_\rho/(w-\rho)\) for zeros in the closed disk centered at
\(21/20+i\Im w\) of radius \(11/60\), plus an error.
If \(C\) is the absolute polynomial-growth constant in
\(|L(s,\chi)|\le C m(1+|s|)\) on \(\Re s\ge1/2\), let
\(A=\log(42(|C|+1))/\log2+1\).
The local lemma bounds the error by
\((1440/11)A\log(m(|\Im w|+2))\), and the sum of the
multiplicities by \(A\log(m(|\Im w|+2))/\log2\).
All these zeros have imaginary part at most \(Q\) in absolute
value, so the box forces \(\Re\rho\le399/400\). Therefore
\(|w-\rho|\ge\delta\), and \(m(|\Im w|+2)\le Q^3\).
In the range \(\Re w\ge11/10\), the Euler-product derivative
bound \(|\ell_\chi(w)|\le 1/(\Re w-1)+C_R\) suffices. One may take
\[
 B_\delta=3A\left(\frac{1440}{11}
               +\frac{1}{\delta\log2}\right)
                  +\frac{10+C_R}{\log2}.
\]


Here is the disk argument underlying those two estimates. Suppose
\(f\) is holomorphic in the disk of radius \(r\) about \(z_0\),
continuous on its closure, \(f(z_0)\ne0\), and
\(|f|\le M|f(z_0)|\) there, \(M\ge1\).
Take every zero in the closed concentric disk of radius \(r/3\),
with its multiplicity, and let \(P\) be the product of the
corresponding monic linear factors. The set is finite by isolated
zeros and compactness. Write \(k=\deg P\) and \(g=f/P\), extending
across the zeros by their local power series. On the outer circle,
\(|P|\ge(2r/3)^k\), whereas \(|P(z_0)|\le(r/3)^k\). The maximum
modulus principle gives
\[
             2^k |g(z)|\le M |g(z_0)|\qquad (|z-z_0|\le r).
\]
At \(z_0\) this implies \(k\le\log M/\log2\). The quotient \(g\)
has no zeros in the disk of radius \(R=r/3\). Take its normalized
holomorphic logarithm \(u=\log(g/g(z_0))\) there. Then \(u(z_0)=0\)
and \(\Re u\le\log M\). For \(a>\log M\), the function
\(u/(2a-u)\) has modulus at most one and vanishes at \(z_0\).
Schwarz's lemma and the triangle inequality give
\[
 |u(z)|\le\frac{2a|z-z_0|}{R-|z-z_0|}\le6a
       \qquad (|z-z_0|\le3R/4).
\]
At any point within \(R/2\) of the center use the Cauchy estimate
on the circle of radius \(R/4\) about that point. Letting
\(a\) decrease to \(\log M\) gives \(|u'|\le24\log M/R\).
Logarithmic differentiation of \(f=Pg\) therefore gives
\[
 \left|\frac{f'(z)}{f(z)}
           -\sum_{|\rho-z_0|\le r/3}\frac{m_\rho}{z-\rho}\right|
     \le72\,\log M/r \qquad(|z-z_0|\le r/6,\ f(z)\ne0).
\]
For \(f=L(\cdot,\chi)\), use \(z_0=21/20+i\Im w\),
\(r=11/20\). The disk lies in \(\Re s\ge1/2\). On its closure
\(1+|s|\le2(|\Im w|+2)\); at the center the Euler-product inverse
bound gives \(1/|L(z_0,\chi)|\le21\). Thus one may use
\(M=42(|C|+1)m(|\Im w|+2)\).
Because \(m(|\Im w|+2)\ge2\), its logarithm satisfies
\(\log M\le A\log(m(|\Im w|+2))\).
The points in the range
\(\sigma\le\Re w<11/10\) satisfy \(|w-z_0|\le r/6\).
For completeness, on \(\Re s>0\) partial summation writes
\(L(s,\chi)=s\int_1^\infty P_\chi(t)t^{-s-1}\,dt\) with
\(P_\chi(t)=\sum_{n\le t}\chi(n)\). Full periods have zero sum,
so \(|P_\chi(t)|\le m\) and \(|L(s,\chi)|\le m|s|/\Re s\).
This supplies the asserted growth on the whole disk. For
\(\Re s>1\), the absolute Möbius series of the inverse Euler
product is bounded by \(\zeta(\Re s)\le1+1/(\Re s-1)\);
at \(z_0\) this is \(21\). Finally the only right-half-plane
bounds used above are on \(\Re s\ge11/10\) and \(\Re s=2\).
Both follow directly by majorizing the logarithmic-derivative
series by the convergent constant
\(\sum_{n\ge2}(\log n)n^{-11/10}\), increasing \(C_R\) if needed.
These computations justify the specializations of the local
expansion, including \(72/r=1440/11\).

For \(y\ge1\) the integrated Perron formula on the line \(2+it\)
gives \(A_\chi(y)\) as the integral of
\[
      F_y(w)=\ell_\chi(w)\frac{y^w}{w(w-1)}
\]
with measure \(dw/(2\pi i)\) upwards. The formula follows by the
absolutely convergent Dirichlet series for \(\ell_\chi\) on this
line and the elementary Mellin inversion of \((u-1)_+\), whose
transform is \(1/[w(w-1)]\). At \(u=1\) the summand is zero from
either side. Shift the finite part of the line between heights
\(-T,T\) to \(\Re w=\sigma\). The box excludes zeros on and
inside the rectangle. The pole of the kernel at \(w=1\)
contributes \(y\ell_\chi(1)\). The other pole, \(w=0\), lies outside the contour.
For the remaining error the two tails on the right are accompanied
by the horizontal segments at \(\pm T\) and the left vertical
segment. Apply this identity at \(y=z+h\) and \(y=z\), and
subtract before estimating the left segments.

On the right \(|\ell_\chi(2+it)|\le1+C_R\) and the integrand
is at most \((1+C_R)y^2/t^2\) for nonzero \(t\).
Since \(y\le2z\), each pair of omitted tails after division
by \(2\pi\) is bounded by \(8(1+C_R)z^2/Q\). On each horizontal
segment the denominator has modulus at least \(T^2\), the length
is at most \(2\), and the logarithmic derivative is bounded by
\(B_\delta L_Q\), since \(T\le Q-1\). Its unnormalized integral
is bounded by \(8B_\delta L_Q z^2/T^2\le16B_\delta L_Q z^2/Q\).

For \(w=\sigma+iu\), integrate the derivative of
\(\exp(tw\log(1+\eta))\), \(0\le t\le1\), to see that
\[
 |(z+h)^w-z^w|\le4\eta |w|z^\sigma .
\]
Thus the factor \(w\) in the kernel cancels in the magnitude of the
difference. Since \(1-\sigma\ge1/800\),
\[
 \int_{-T}^T\frac{du}{|\sigma-1+iu|}
 \le\int_{-T}^T\frac{2\,du}{1/800+|u|}
 =4\log(1+800T)\le40\log Q .
\]
The left difference before division by \(2\pi\) is at most
\(160 B_\delta L_Q^2\eta z^\sigma\). There are two pairs of
right tails and four horizontal segments. Using
\(L_Q\ge\log2\ge1/2\) and \(2\pi\ge1\), their bounds and
the vertical bound give \eqref{eq:box-contour} with
\(C_\delta=160B_\delta+32(1+C_R)\).

The independent error at the finite cutoff can be read directly
from the finite sum defining \(A_\chi\):
\[
 A_\chi(z+h)-A_\chi(z)-h\sum_{n\le z}a_n/n
     =\sum_{z<n\le z+h}a_n((z+h)/n-1).
\]
There are at most \(h+1\) summands, \(|a_n|\le\log(z+h)\),
and the factor is between zero and \(h/z\). Subtracting the
contour expression and dividing by \(h\) proves
\begin{equation}
 |R_\chi(z)|\le
 \frac{C_\delta L_Q\{L_Q\eta z^\sigma+z^2/Q\}
              +h(h+1)\log(z+h)/z}{h}.
 \label{eq:box-tail}
\end{equation}
This estimate applies to a general nonprincipal character obeying
the box and modulus conditions.

\subsection{Uniformity across a real interval}

Fix \(0<\epsilon\le1/4,\ \tau>0,\ K\ge1\) and take
\[
    Y^{\epsilon/6}\le \log Q\le Y^{\epsilon/4},\qquad
    x=Y^K,\quad z=x^{400},\quad \eta=x^{-2},\quad
    h=x^{398},\quad \delta=\epsilon/(1600K).
\]
In particular \(0<\delta\le1/800\).
Multiplication of \eqref{eq:box-tail} by \(x\) gives exactly
\[
 x|R_\chi(z)|\le C_\delta(\log Q)^2Y^{\epsilon/4}
   +C_\delta(\log Q)x^{403}/Q
   +\frac{(x^{398}+1)\log(x^{400}+x^{398})}{x^{399}} .
\]
Indeed \(z^\sigma=x^{399}Y^{\epsilon/4}\).
The first term is at most \(C_\delta Y^{3\epsilon/4}\), hence at
most \(Y^\epsilon/3\) for large \(Y\). Since \(\log Q\le x\)
and \(Q\ge\exp(Y^{\epsilon/6})\), the second is bounded by
\(C_\delta Y^{404K}\exp(-Y^{\epsilon/6})\), which is eventually
at most \(1/3\). For \(x\ge3\), \(x^{400}+x^{398}\le x^{401}\)
and \(x^{398}+1\le2x^{398}\). The last term is at most
\(802\log x/x\le802\), eventually at most \(Y^\epsilon/3\).
Increasing the common threshold accordingly gives
\[
                      x|R_\chi(x^{400})|\le Y^\epsilon .
\]
Now let \(q\) be the prime index and suppose
\(\chi(n)=\kappa(n)(\frac{n}{q})\) for every positive integer
\(n\). This is the primitive twist selected at the counting stage.
If throughout \(Y\le y\le Y^K\) the finite sum at
\(z=(Y^K)^{400}\) has absolute value at most
\(y^{1/2+\tau}/2\), then \eqref{eq:box-finite-split} yields
\[
 |G_\chi(y)|\le \tfrac12y^{1/2+\tau}+y|R_\chi(z)|
                    \le y^{1/2+\tau}.
\]
The last step uses \(y\le x\), and, for \(Y\ge16\),
\(Y^\epsilon\le Y^{1/4}\le Y^{1/2}/2\le y^{1/2+\tau}/2\).
The single choice of truncation applies even when the largest
finite sum is attained at a different \(y\) for each prime.

\subsection{The constant four in the logarithmic derivative bound}

The immediate bound in \texttt{box\_of\_S} is
\(|L'/L(1,\chi)|\le C\log Y\). Indeed apply the preceding tail
estimate at \(K=1\) to get \(|R_\chi(Y^{400})|\le1\), then use
\(\sum_{n\le t}\Lambda(n)/n\le C_0\log t\) with \(t=Y^{400}\). The constant
\(400C_0+1\) suffices for \(Y\ge3\).
To obtain \(4\), the source uses the sharper
\(\sum_{n\le t}\Lambda(n)/n\le\log t+2\) and changes the
truncation. For the latter Mertens estimate take \(N=\lfloor t\rfloor\).
The identity \(\sum_{d\mid n}\Lambda(d)=\log n\) yields
\[
 N\sum_{n\le N}\Lambda(n)/n
 \le\sum_{n\le N}\Lambda(n)\lfloor N/n\rfloor+\psi(N)
 =\log(N!)+\psi(N).
\]
The increasing-function integral bound gives
\(\log(N!)\le N\log N-N+1+\log N\).
The source's explicit elementary Chebyshev estimate
\(\psi(N)+1+\log N\le3N\) gives the claimed result after
division by \(N\).

Write \(l=\log Q,\ s=l^{1/4},\ x=s^{10}=l^{5/2}\). Now use
\(\delta=(2/5)/1600\) in \eqref{eq:box-tail}, with the same
\(\eta=x^{-2}\), \(z=x^{400}\). The numerical bound becomes
\[
 x|R_\chi(z)|\le C_\delta l^2 s+
                 C_\delta l x^{403}/Q+802 .
\]
If \(l\) is sufficiently large then \(l\ge81\), so \(x\ge3\),
and \(l^{1006}\le e^l=Q\). The first term without its constant
is \(s^9\le x\). The second satisfies \(lx^{403}=x l^{1006}\le xQ\).
Thus \(|R_\chi(z)|\le 2C_\delta+802\). Since
\(\log z=1000\log\log Q\), the Mertens estimate gives an absolute
constant \(C'\) such that
\[
       |L'(1,\chi)/L(1,\chi)|\le1000\log\log Q+C'.
\]
For the pair windows it suffices that \(0<\epsilon<1/52\) and,
for fixed \(A>0\),
\(
 Y^{\epsilon/6}\le\log Q\le A Y^{\epsilon/5}(\log Y)^2.
\)
The lower bound ensures the preceding threshold in \(l\) for
large \(Y\). Writing \(L=\log Y\), the upper bound gives
\[
 1000\log\log Q+C'\le200\epsilon L+2000\log L+1000\log A+C'.
\]
The first summand is at most \((50/13)L\), the second is at most
\(L/13\) eventually, and the two fixed constants together are at
most \(L/13\) eventually. Thus the expression is at most
\(4L\le4\log y\) for every \(y\in[Y,Y^K]\). In the actual
pair selection \(A=3\) suffices. The lower threshold depends
on the fixed parameters only. The constant four and the
prime-sum bound therefore hold throughout the same interval.

\paragraph{Source correspondence.}
The deductions are \texttt{BoxOfS.box\_of\_S}, \texttt{TailKExp.tailK},
\texttt{kappaC.kappaC} and \texttt{Kappa4Loglog.kappa4}, under
\srcid{Development.AutomorphicArithmetic.MeanSquareBounds}.
The rectangle, pointwise logarithmic-derivative bound, finite
difference and numerical substitutions are \texttt{TailContour},
\texttt{TailBox}, \texttt{TailDiff}, \texttt{TailKernel} and
\texttt{TailNum} under
\srcid{Formalism.GL3MeanSquares.GenusMeanAsymptotics.JacobiCharacterSums}.
The local zero expansion is proved above from its disk factorization.
The polynomial growth and Euler-product bounds proved above correspond to \texttt{LandauLemmaLfunctionZeros},
\texttt{LGrowth} and \texttt{LRight}. The exact elementary Mertens
bound and its explicit Chebyshev input occur in \texttt{mertensUpper}
under \srcid{Development.AutomorphicArithmetic.DedekindLFunctions.LandauTheorems}.
